% questions.tex —— 1958 年·湖北卷 · 题目内容片段（无答案）
% 由 tools/generate_projects.py 依来源总账自 1958/1958_湖北_数学真题_试题.md 生成，勿手改（改源 MD 后重跑）


\begin{question}
\gkti{1}设有数列 $1+2x+3x^2+4x^3+\cdots\cdots$

（1） 写出这个数列的第 $n$ 项；

（2） 写出这个数列的前 $n$ 项和的公式；

（3） 问这个和的公式当 $x=1$ 时能否成立？如能成立，说明理由，如不能成立，怎样求和？

\end{question}
\begin{question}
\gkti{2}设 $A$、$B$、$C$ 是 $\triangle ABC$ 的三个内角，求证：

\nopagebreak[4]
\[
\tan A+\tan B+\tan C=\tan A\cdot\tan B\cdot\tan C
\]

\end{question}
\begin{question}
\gkti{3}解方程：

\nopagebreak[4]
\[
\left(\log_x\sqrt5\right)^2+\log_x 5\sqrt5+\frac54=0
\]

\end{question}
\begin{question}
\gkti{4}在 $\triangle ABC$ 中，$AD$ 是 $BC$ 的中线，而且 $\angle B$ 和 $\angle CAD$ 互为余角，问 $\triangle ABC$ 是怎样的三角形？

\end{question}
\begin{question}
\gkti{5}在直角三角形所在平面外取一点 $P$，而且 $P$ 点和三角形三个顶点的距离都是相等的，证明：连接 $P$ 点和三角形斜边中点的直线垂直于这个三角形所在的平面。

\begin{gktip}
\textbf{源注}：原卷第 5 题之后有居中标题「补充题（非统考题）」，其下三题接续正式题编号为 $6$–$8$；按本库约定以「附加题」挂于正式题之后，编号照原卷不变，不设小节标题。
\end{gktip}

\textbf{附加题 6．}比较 $\left(\dfrac{21}{20}\right)^{100}$ 与 $100$ 的大小。

\Needspace{5\baselineskip}
\textbf{附加题 7．}已知 $a^2+b^2=7ab$，求证：

\nopagebreak[4]
\[
\lg\left(\frac{a+b}{3}\right)=\frac12(\lg a+\lg b)
\]

\textbf{附加题 8．}已知方程式 $x^3+px^2+qx+r=0$ 的三根 $a$、$b$、$c$ 成调和级数，求 $a$、$b$、$c$。

\end{question}